\section{Discussion}
\label{sec:discussion}

\subsection{When is self-training UDA useful?}
The asymmetric gain finding translates into a practical decision rule.
Before committing engineering effort to a self-training UDA pipeline,
practitioners should first measure the source-only target $\hmean$
under a post-process sweep. If the swept-best $\hmean$ is bounded
below $\sim$$0.01$ (no detectable target signal at any threshold), the
UDA pipeline is the only viable path to non-trivial detection and the
gain can be substantial -- $+6.26$~pp in our $\stoh$ direction. If
the swept-best $\hmean$ already exceeds the noise floor, threshold
sweeping and modest backbone upgrades are likely to recover the same
gain at a fraction of the engineering cost; the additional plumbing
of an EMA teacher and a domain head is unlikely to repay its complexity
cost. This regime distinction is, in our view, the most actionable
output of the present study.

\subsection{Why does long training collapse?}
We attribute the long-run collapse to a teacher--student drift loop
that is intrinsic to confidence-thresholded self-training without
source-target consistency regularisation. Once the EMA teacher tilts
towards over-confident predictions (high precision, low recall on the
target validation set, observable in our per-epoch metrics from
epoch~$9$ onwards on $\hsource$), the binarisation step in
Eq.~\ref{eq:keep} retains only the highest-confidence regions as
pseudo-ground-truth. The student is then trained to reproduce these
increasingly conservative pseudo-labels, which the EMA step in
Eq.~\ref{eq:ema} folds back into the teacher with a $1{-}\alpha
=10^{-3}$ coefficient. Over $\sim$$10$ epochs the loop tightens,
recall collapses, and the model converges to predicting nothing -- the
same failure mode as the cold-start case, reached this time from above
rather than from below.

Three mechanisms could in principle break the loop. (i) A consistency
loss between weak and strong augmentations of the same target image
would give the student a non-pseudo signal on the target domain.
(ii) A confidence-aware loss weight on each pseudo-pixel, rather than
the coarse per-image keep flag of Eq.~\ref{eq:keep}, would let
low-confidence regions contribute proportionally rather than be
dropped. (iii) A teacher reset or rollback policy could escape the
drift trajectory once collapse is detected. We leave a systematic
comparison of these mitigations to future work; in the present paper
the $5$-epoch budget side-steps the problem at no measurable cost to
final $\hmean$.

\subsection{Limitations}
The present study has three limitations that we name explicitly
rather than minimise. First, the backbone is ResNet-18 throughout;
MPS memory constraints on our single-device setting made longer
ResNet-50 schedules impractical, and we have not verified whether the
asymmetric gain phenomenon persists or attenuates with a stronger
oracle. Second, the domain head is image-level only; a multi-scale or
region-level DANN may close more of the $\hsource$ gap, where our
results suggest weak feature alignment rather than weak detection
capability is the binding constraint. Third, we report single-seed
numbers because the per-run compute budget (5--10 min/epoch on MPS) is
small but multi-run sweeps were still time-consuming; our direct
re-runs of v3 give a $\sim$$0.4$~pp $\hmean$ run-to-run variance on
validation ($0.0596$ vs.\ $0.0633$), which is comparable to the
$\hsource$ UDA gain itself, and we have flagged this explicitly in
the headline interpretation.

\subsection{Future work}
Three extensions are natural. (i) Replace the ResNet-18 + image-DANN
combination with a ResNet-50 + multi-scale DANN, and re-measure both
the absolute $\hmean$ and the magnitude of the asymmetric gap. We
expect the cold-start regime ($\stoh$) to widen further -- a stronger
source encoder makes the cold-start fix even more powerful -- while
the $\hsource$ regime may move closer to its calibration ceiling.
(ii) Add a MIC-style~\cite{hoyer2023mic} masked-image consistency loss
to the pseudo-label branch, replacing the keep-or-drop binary filter
of Eq.~\ref{eq:keep} with a soft consistency objective that we
hypothesise would break the drift collapse and unlock longer training.
(iii) Once Mongolian recognition annotations become available,
extend the pipeline to end-to-end spotting; the same cross-domain
asymmetry is likely to appear at the recognition layer with a
different (and possibly more aggressive) cold-start failure mode.
